\subsection*{最初に必要な用語}

\addcontentsline{toc}{subsection}{最初に必要な用語}

\begin{center}
\small
\par\smallskip\noindent\refstepcounter{table}\label{tab:ch14-01}表 \thetable: 最初に必要な用語の整理\par\smallskip
表 \thetable は、最初に必要な用語の整理を比較・確認しやすい形で整理したものである。\par\smallskip
\begin{longtable}{p{0.24\linewidth}p{0.71\linewidth}}
\toprule
用語 & 初学者向けの説明 \\
\midrule
\endfirsthead
\toprule
用語 & 初学者向けの説明 \\
\midrule
\endhead
\term{専門職実務}{Professional Practice} & 社会的責任を伴う仕事を、一定の標準・倫理・技能に従って行うこと。 \\
\midrule
専門性 & 専門分野の知識だけでなく、責任ある判断、継続学習、職業倫理を含む態度である。 \\
\midrule
倫理規程 & 専門家が守るべき価値観と行動原則を明文化したもの。 \\
\midrule
専門職行動規範 & 専門家として実際に避けるべき行為、行うべき行為を示す規範である。 \\
\midrule
専門職団体 & 専門分野の知識体系、標準、資格、倫理、教育、交流を支える団体である。 \\
\midrule
標準 & 成果物やプロセスが最低限満たすべき性質や手順を定めた合意文書である。 \\
\midrule
認定 & 教育機関やプログラムが一定の基準を満たすことを第三者が認めること。 \\
\midrule
資格証明 & 個人が特定分野の知識・能力を持つことを試験などで確認すること。 \\
\midrule
免許 & 特定の業務を行う権限を、政府または法定機関が個人に与えること。 \\
\midrule
知的財産 & 発明、著作物、設計、商標、営業秘密など、知的活動から生まれる財産的価値である。 \\
\midrule
利害関係者 & ソフトウェアの開発・運用結果に影響を与える、または影響を受ける人や組織である。 \\
\midrule
トレードオフ & 一方を良くすると他方が悪くなる関係の中で、目的に照らして釣り合いを取る判断である。 \\
\midrule
集団力学 & チームや組織の中で、人の行動、協力、衝突、意思決定がどのように生じるかを扱う考え方である。 \\
\midrule
包摂性 & 性別、文化、言語、地域、障害、経験などの違いを前提に、人が参加しやすい状態を作ること。 \\
\midrule
\bottomrule
\end{longtable}
\normalsize
\end{center}
